\documentclass[10pt,a4paper]{amsart}
\usepackage{amsmath,amssymb,amsthm,graphicx}
\usepackage{verbatim}
\newenvironment{reference}{\comment}{\endcomment}
%\newenvironment{reference}{}{}
\newenvironment{slogan}{\comment}{\endcomment}
\newenvironment{history}{\comment}{\endcomment}

% For commutative diagrams we use Xy-pic
\usepackage[all]{xy}

% We use 2cell for 2-commutative diagrams.
\xyoption{2cell}
% Theorem environments.
%
\theoremstyle{plain}
\newtheorem{theorem}[subsection]{Theorem}
\newtheorem{proposition}[subsection]{Proposition}
\newtheorem{lemma}[subsection]{Lemma}

\theoremstyle{definition}
\newtheorem{definition}[subsection]{Definition}
\newtheorem{example}[subsection]{Example}
\newtheorem{exercise}[subsection]{Exercise}
\newtheorem{situation}[subsection]{Situation}

\theoremstyle{remark}
\newtheorem{remark}[subsection]{Remark}
\newtheorem{remarks}[subsection]{Remarks}

\numberwithin{equation}{subsection}

% Macros
%
\def\lim{\mathop{\mathrm{lim}}\nolimits}
\def\colim{\mathop{\mathrm{colim}}\nolimits}
\def\Spec{\mathop{\mathrm{Spec}}}
\def\Hom{\mathop{\mathrm{Hom}}\nolimits}
\def\Ext{\mathop{\mathrm{Ext}}\nolimits}
\def\SheafHom{\mathop{\mathcal{H}\!\mathit{om}}\nolimits}
\def\SheafExt{\mathop{\mathcal{E}\!\mathit{xt}}\nolimits}
\def\Sch{\mathit{Sch}}
\def\Mor{\mathop{\mathrm{Mor}}\nolimits}
\def\Ob{\mathop{\mathrm{Ob}}\nolimits}
\def\Sh{\mathop{\mathit{Sh}}\nolimits}
\def\NL{\mathop{N\!L}\nolimits}
\def\CH{\mathop{\mathrm{CH}}\nolimits}
\def\proetale{{pro\text{-}\acute{e}tale}}
\def\etale{{\acute{e}tale}}
\def\QCoh{\mathit{QCoh}}
\def\Ker{\mathop{\mathrm{Ker}}}
\def\Im{\mathop{\mathrm{Im}}}
\def\Coker{\mathop{\mathrm{Coker}}}
\def\Coim{\mathop{\mathrm{Coim}}}

% Boxtimes
%
\DeclareMathSymbol{\boxtimes}{\mathbin}{AMSa}{"02}

%
% Macros for moduli stacks/spaces
%
\def\QCohstack{\mathcal{QC}\!\mathit{oh}}
\def\Cohstack{\mathcal{C}\!\mathit{oh}}
\def\Spacesstack{\mathcal{S}\!\mathit{paces}}
\def\Quotfunctor{\mathrm{Quot}}
\def\Hilbfunctor{\mathrm{Hilb}}
\def\Curvesstack{\mathcal{C}\!\mathit{urves}}
\def\Polarizedstack{\mathcal{P}\!\mathit{olarized}}
\def\Complexesstack{\mathcal{C}\!\mathit{omplexes}}
% \Pic is the operator that assigns to X its picard group, usage \Pic(X)
% \Picardstack_{X/B} denotes the Picard stack of X over B
% \Picardfunctor_{X/B} denotes the Picard functor of X over B
\def\Pic{\mathop{\mathrm{Pic}}\nolimits}
\def\Picardstack{\mathcal{P}\!\mathit{ic}}
\def\Picardfunctor{\mathrm{Pic}}
\def\Deformationcategory{\mathcal{D}\!\mathit{ef}}
\usepackage{cleveref}
\crefname{originallemma}{original lemma}{original lemmas}
\Crefname{originallemma}{Original lemma}{Original lemmas}
\makeatletter
\newlabel{descent-lemma-compare-etale-zariski-flat}{{0GNA}{source}{}{sourcepointer-1}{}}
\newlabel{descent-lemma-compare-etale-zariski-flat@cref}{{[originallemma][0][]0GNA}{[1][0][]source}{}{}{}}
\makeatother
\crefname{originalequation}{original equation}{original equations}
\Crefname{originalequation}{Original equation}{Original equations}
\makeatletter
\newlabel{equation-epsilon}{{08H7}{source}{}{sourcepointer-2}{}}
\newlabel{equation-epsilon@cref}{{[originalequation][0][]08H7}{[1][0][]source}{}{}{}}
\makeatother
\crefname{originallemma}{original lemma}{original lemmas}
\Crefname{originallemma}{Original lemma}{Original lemmas}
\makeatletter
\newlabel{etale-cohomology-lemma-describe-etale-local-ring}{{04HX}{source}{}{sourcepointer-3}{}}
\newlabel{etale-cohomology-lemma-describe-etale-local-ring@cref}{{[originallemma][0][]04HX}{[1][0][]source}{}{}{}}
\makeatother
\crefname{originalreference}{original reference}{original references}
\Crefname{originalreference}{Original reference}{Original references}
\makeatletter
\newlabel{etale-cohomology-remarks-enough-points}{{040S}{source}{}{sourcepointer-4}{}}
\newlabel{etale-cohomology-remarks-enough-points@cref}{{[originalreference][0][]040S}{[1][0][]source}{}{}{}}
\makeatother
\crefname{originallemma}{original lemma}{original lemmas}
\Crefname{originallemma}{Original lemma}{Original lemmas}
\makeatletter
\newlabel{more-algebra-lemma-dumb-properties-henselization}{{07QM}{source}{}{sourcepointer-5}{}}
\newlabel{more-algebra-lemma-dumb-properties-henselization@cref}{{[originallemma][0][]07QM}{[1][0][]source}{}{}{}}
\makeatother
\crefname{originallemma}{original lemma}{original lemmas}
\Crefname{originallemma}{Original lemma}{Original lemmas}
\makeatletter
\newlabel{sites-lemma-point-morphism-sites}{{05V1}{source}{}{sourcepointer-6}{}}
\newlabel{sites-lemma-point-morphism-sites@cref}{{[originallemma][0][]05V1}{[1][0][]source}{}{}{}}
\makeatother
\crefname{originallemma}{original lemma}{original lemmas}
\Crefname{originallemma}{Original lemma}{Original lemmas}
\makeatletter
\newlabel{sites-modules-lemma-check-flat-stalks}{{05VC}{source}{}{sourcepointer-7}{}}
\newlabel{sites-modules-lemma-check-flat-stalks@cref}{{[originallemma][0][]05VC}{[1][0][]source}{}{}{}}
\makeatother
\crefname{originallemma}{original lemma}{original lemmas}
\Crefname{originallemma}{Original lemma}{Original lemmas}
\makeatletter
\newlabel{sites-modules-lemma-flat-pullback-exact}{{04JC}{source}{}{sourcepointer-8}{}}
\newlabel{sites-modules-lemma-flat-pullback-exact@cref}{{[originallemma][0][]04JC}{[1][0][]source}{}{}{}}
\makeatother
\crefname{originallemma}{original lemma}{original lemmas}
\Crefname{originallemma}{Original lemma}{Original lemmas}
\makeatletter
\newlabel{sites-modules-lemma-pullback-stalk}{{05V5}{source}{}{sourcepointer-9}{}}
\newlabel{sites-modules-lemma-pullback-stalk@cref}{{[originallemma][0][]05V5}{[1][0][]source}{}{}{}}
\makeatother
\title{The morphism from the Zariski to the etale site is flat and conservative}
\author{The Stacks Project Authors}
\date{stacks-project:spaces-perfect.tex:local-snapshot-20261012; GNU FDL 1.2 or later}
\begin{document}
\maketitle
\noindent\textit{Editorial note.} All delivered mathematical text is the author's own, copied from the pinned Stacks Project chapter file spaces-perfect.tex; no statement or proof step was rewritten, reordered or removed.\par\medskip
\noindent\textit{Editorial note.} The counted proof is the author's complete argument: the proof checks flatness on stalks through the strict henselization description of the etale local ring, derives exactness of epsilon\textasciicircum{}* from flatness, and proves faithfulness by tensoring at each stalk and applying faithful flatness\par\medskip
\noindent\textit{Editorial note.} The label lemma-epsilon-flat names the counted statement, and the proof environment immediately following it in the source is delivered verbatim as its proof.\par\medskip
\noindent\textit{Editorial note.} The counted statement is carried by the author's own lemma environment in the chapter source.\par\medskip
\noindent\textit{Allowed standard prerequisites (editorial).}\begin{itemize}
\item the localization morphism of ringed sites epsilon from the small etale site of X to its Zariski site, and the stalk formula for epsilon\textasciicircum{}*F at a geometric point
\item flatness of a morphism of ringed sites checked on stalks (Modules on Sites, Lemma sites-modules-lemma-check-flat-stalks), the point morphism of sites (Sites, Lemma sites-lemma-point-morphism-sites) and enough points (Etale Cohomology, Remarks etale-cohomology-remarks-enough-points)
\item the etale local ring O\_\{etale, xbar\} is the strict henselization of O\_\{X,x\} (Etale Cohomology, Lemma etale-cohomology-lemma-describe-etale-local-ring)
\item faithful flatness of the strict henselization over the local ring (More on Algebra, Lemma more-algebra-lemma-dumb-properties-henselization)
\item flat pullback of modules is exact (Modules on Sites, Lemma sites-modules-lemma-flat-pullback-exact)
\end{itemize}\medskip
\ifdefined\hypertarget\hypertarget{sourcepointer-1}{}\fi
\noindent\textit{Source reference.} Original lemma 0GNA: sources/stacks-chow/tags.txt, source lines 20785--20785. Prerequisite pointer invoked at proof line(s) 340; it names a standard fact of a sibling chapter that the counted argument uses, and every step of the counted proof built on it is retained verbatim in the delivered body.\par\medskip
\ifdefined\hypertarget\hypertarget{sourcepointer-2}{}\fi
\noindent\textit{Source reference.} Original equation 08H7: sources/stacks-spaces-perfect/pinned-original-source.tex, source lines 320--320. Prerequisite pointer invoked at proof line(s) 331; it names a standard fact of the same chapter that the counted argument uses, and every step of the counted proof built on it is retained verbatim in the delivered body.\par\medskip
\ifdefined\hypertarget\hypertarget{sourcepointer-3}{}\fi
\noindent\textit{Source reference.} Original lemma 04HX: sources/stacks-chow/tags.txt, source lines 5883--5883. Prerequisite pointer invoked at proof line(s) 354; it names a standard fact of a sibling chapter that the counted argument uses, and every step of the counted proof built on it is retained verbatim in the delivered body.\par\medskip
\ifdefined\hypertarget\hypertarget{sourcepointer-4}{}\fi
\noindent\textit{Source reference.} Original reference 040S: sources/stacks-chow/tags.txt, source lines 5267--5267. Prerequisite pointer invoked at proof line(s) 352; it names a standard fact of a sibling chapter that the counted argument uses, and every step of the counted proof built on it is retained verbatim in the delivered body.\par\medskip
\ifdefined\hypertarget\hypertarget{sourcepointer-5}{}\fi
\noindent\textit{Source reference.} Original lemma 07QM: sources/stacks-chow/tags.txt, source lines 9837--9837. Prerequisite pointer invoked at proof line(s) 359; it names a standard fact of a sibling chapter that the counted argument uses, and every step of the counted proof built on it is retained verbatim in the delivered body.\par\medskip
\ifdefined\hypertarget\hypertarget{sourcepointer-6}{}\fi
\noindent\textit{Source reference.} Original lemma 05V1: sources/stacks-chow/tags.txt, source lines 7535--7535. Prerequisite pointer invoked at proof line(s) 349; it names a standard fact of a sibling chapter that the counted argument uses, and every step of the counted proof built on it is retained verbatim in the delivered body.\par\medskip
\ifdefined\hypertarget\hypertarget{sourcepointer-7}{}\fi
\noindent\textit{Source reference.} Original lemma 05VC: sources/stacks-chow/tags.txt, source lines 7546--7546. Prerequisite pointer invoked at proof line(s) 347; it names a standard fact of a sibling chapter that the counted argument uses, and every step of the counted proof built on it is retained verbatim in the delivered body.\par\medskip
\ifdefined\hypertarget\hypertarget{sourcepointer-8}{}\fi
\noindent\textit{Source reference.} Original lemma 04JC: sources/stacks-chow/tags.txt, source lines 5933--5933. Prerequisite pointer invoked at proof line(s) 364; it names a standard fact of a sibling chapter that the counted argument uses, and every step of the counted proof built on it is retained verbatim in the delivered body.\par\medskip
\ifdefined\hypertarget\hypertarget{sourcepointer-9}{}\fi
\noindent\textit{Source reference.} Original lemma 05V5: sources/stacks-chow/tags.txt, source lines 7539--7539. Prerequisite pointer invoked at proof line(s) 373; it names a standard fact of a sibling chapter that the counted argument uses, and every step of the counted proof built on it is retained verbatim in the delivered body.\par\medskip
\begin{lemma}
\label{lemma-epsilon-flat}
The morphism $\epsilon$ of (\ref{equation-epsilon})
is a flat morphism of ringed sites. In particular the functor
$\epsilon^* : \textit{Mod}(\mathcal{O}_X) \to
\textit{Mod}(\mathcal{O}_\etale)$ is exact.
Moreover, if $\epsilon^*\mathcal{F} = 0$, then $\mathcal{F} = 0$.
\end{lemma}

\begin{proof}
The flatness of the morphism $\epsilon$ is
Descent, Lemma \ref{descent-lemma-compare-etale-zariski-flat}.
Here is another proof. We have to show that
$\mathcal{O}_\etale$ is a flat $\epsilon^{-1}\mathcal{O}_X$-module.
To do this it suffices to check
$\mathcal{O}_{X, x} \to \mathcal{O}_{\etale, \overline{x}}$
is flat for any geometric point $\overline{x}$ of $X$, see
Modules on Sites, Lemma
\ref{sites-modules-lemma-check-flat-stalks},
Sites, Lemma
\ref{sites-lemma-point-morphism-sites},
and
\'Etale Cohomology, Remarks
\ref{etale-cohomology-remarks-enough-points}.
By \'Etale Cohomology, Lemma
\ref{etale-cohomology-lemma-describe-etale-local-ring}
we see that $\mathcal{O}_{\etale, \overline{x}}$ is the
strict henselization of $\mathcal{O}_{X, x}$. Thus
$\mathcal{O}_{X, x} \to \mathcal{O}_{\etale, \overline{x}}$
is faithfully flat by More on Algebra,
Lemma \ref{more-algebra-lemma-dumb-properties-henselization}.

\medskip\noindent
The exactness of $\epsilon^*$ follows from the flatness
of $\epsilon$ by
Modules on Sites, Lemma \ref{sites-modules-lemma-flat-pullback-exact}.

\medskip\noindent
Let $\mathcal{F}$ be an $\mathcal{O}_X$-module. If
$\epsilon^*\mathcal{F} = 0$, then with notation as above
$$
0 = \epsilon^*\mathcal{F}_{\overline{x}} =
\mathcal{F}_x \otimes_{\mathcal{O}_{X, x}} \mathcal{O}_{\etale, \overline{x}}
$$
(Modules on Sites, Lemma \ref{sites-modules-lemma-pullback-stalk})
for all geometric points $\overline{x}$. By faithful flatness of
$\mathcal{O}_{X, x} \to \mathcal{O}_{\etale, \overline{x}}$
we conclude $\mathcal{F}_x = 0$ for all $x \in X$.
\end{proof}

\end{document}
